\documentclass[11pt,a4paper]{article}

\usepackage[margin=2.5cm]{geometry}
\usepackage{booktabs}
\usepackage{array}
\usepackage{tabularx}
\usepackage{parskip}
\usepackage{microtype}
\usepackage{enumitem}
\usepackage{fancyhdr}
\usepackage{titlesec}
\usepackage{hyperref}
\usepackage{caption}

\hypersetup{
  colorlinks=false,
  hidelinks
}

\captionsetup{
  font=small,
  labelfont=bf,
  skip=4pt
}

\setlist[itemize]{nosep, topsep=2pt, partopsep=0pt, leftmargin=*}
\setlist[enumerate]{nosep, topsep=2pt, partopsep=0pt, leftmargin=*}

\titleformat{\section}{\normalfont\large\bfseries}{\thesection.}{0.5em}{}
\titleformat{\subsection}{\normalfont\normalsize\bfseries}{\thesubsection.}{0.5em}{}
\titlespacing*{\section}{0pt}{10pt}{4pt}
\titlespacing*{\subsection}{0pt}{6pt}{2pt}

\pagestyle{fancy}
\fancyhf{}
\renewcommand{\headrulewidth}{0.4pt}
\fancyhead[L]{\small BSE4104 --- Group B}
\fancyhead[R]{\small Week 7 Progress Report}
\fancyfoot[C]{\thepage}

\newcolumntype{L}[1]{>{\raggedright\arraybackslash}p{#1}}
\newcolumntype{C}[1]{>{\centering\arraybackslash}p{#1}}
\newcolumntype{R}[1]{>{\raggedleft\arraybackslash}p{#1}}

% ────────────────────────────────────────────────────────────────────────────
\begin{document}

\begin{center}
  {\LARGE\bfseries Centenary Bank Helpdesk Agent}\\[6pt]
  {\large\bfseries Week 7 Progress Report}\\[4pt]
  {\large Evaluation, Observability and Guardrails}\\[10pt]
  {\normalsize
    BSE4104 Emerging Trends in Software Engineering\\
    Group B \quad|\quad October 2026
  }
\end{center}

\vspace{4pt}
\noindent\rule{\textwidth}{0.4pt}
\vspace{4pt}

% ────────────────────────────────────────────────────────────────────────────
\section{Introduction}

Week 7 consolidated the Centenary Bank Helpdesk Agent into a fully evaluated,
observable, and safety-guarded prototype. Building on the persistent-memory
agentic loop completed in Week 6, this week delivered four interconnected
deliverables: a 30-scenario evaluation set, a structured observability and
tracing bundle, a failure catalogue, and a guardrail implementation module.
Together these demonstrate that the agent can be evaluated systematically
against a defined test set, that every decision is fully auditable, and that
explicit, independently testable safety controls are in place.

% ────────────────────────────────────────────────────────────────────────────
\section{Evaluation Framework}

\subsection{30-Scenario Evaluation Set}

A formal evaluation set was designed to exercise the agent across six scenario
categories, with five cases in each. Each scenario specifies the user request,
the expected agent behaviour, the correct tool path, and any safety flags that
should be recorded. The categories are summarised in Table~\ref{tab:eval}.

\begin{table}[h]
  \centering
  \caption{Evaluation categories and their focus areas}
  \label{tab:eval}
  \small
  \begin{tabular}{@{} L{3.0cm} C{1.0cm} L{8.0cm} @{}}
    \toprule
    \textbf{Category} & \textbf{Cases} & \textbf{Focus} \\
    \midrule
    Normal            & 5 & Branch hours, account opening, card loss, transfer fees, ticket creation \\
    Edge              & 5 & Empty input, ambiguous request, multi-intent, noisy text, long message \\
    Incorrect info    & 5 & False policy claims, wrong ticket status, fake balance, stale rule, nonexistent source \\
    Adversarial       & 5 & Prompt injection, reasoning extraction, secret exfiltration, malicious retrieval, fabrication \\
    Tool failure      & 5 & Empty retrieval, timeout, malformed output, service error, partial data \\
    Unauthorised      & 5 & Cross-customer access, account change, funds transfer, loan approval, credential request \\
    \bottomrule
  \end{tabular}
\end{table}

\subsection{Measurable Metrics}

Six metrics were defined for every scenario: task completion, groundedness,
tool selection accuracy, instruction following, latency, and safety violations.
The evaluation harness (\texttt{week7\_evaluation\_harness.py}) executes the
real Week 6 controller against all 30 scenarios using a scripted planner and
records structured trace events for each run. Suggested latency budgets were
set at 3{,}000\,ms for direct answers, 8{,}000\,ms for tool-using scenarios,
and 12{,}000\,ms for failure-recovery paths.

% ────────────────────────────────────────────────────────────────────────────
\section{Observability and Tracing}

Every evaluation run emits a structured trace stored in
\texttt{evidence/traces/week7/week7-observability-evidence.json}. Each record
captures the fields defined in the Week 7 trace specification:

\begin{itemize}
  \item \textbf{Prompt fingerprint.} Template identifier and SHA-256 hash of
        the rendered prompt, so that any prompt change between runs is
        detectable by comparing hashes.
  \item \textbf{Retrieval context.} The query issued to the knowledge base,
        the IDs of retrieved passages, and the relevance scores returned.
  \item \textbf{Tool trace.} The tool name selected, redacted arguments
        (identity fields removed before logging), the tool status, and the
        error type when the call failed.
  \item \textbf{Latency.} Per-iteration planning latency and total end-to-end
        latency in milliseconds.
  \item \textbf{Terminal outcome.} The controller's final state label, such as
        \texttt{TICKET\_CREATED} or \texttt{OUT\_OF\_SCOPE\_HANDOFF}.
  \item \textbf{Safety flags.} Structured labels emitted when a guardrail
        fires or a policy boundary is crossed during the run.
\end{itemize}

The trace design satisfies the requirement that any run can be fully
reconstructed offline without relying on console output, enabling retrospective
audits and regression testing.

% ────────────────────────────────────────────────────────────────────────────
\section{Failure Catalogue}

Five controlled failure conditions were identified, exercised under the
evaluation harness, and remediated. Each entry in the catalogue records the
failure type, its root cause, the remediation applied, and a reference to the
re-test trace evidence. All five cases re-tested as \textit{Pass}.

\begin{table}[h]
  \centering
  \caption{Failure catalogue --- five cases, all remediated and re-tested}
  \label{tab:failures}
  \small
  \begin{tabular}{@{} C{0.8cm} L{3.0cm} L{4.2cm} L{4.0cm} @{}}
    \toprule
    \textbf{ID} & \textbf{Failure type} & \textbf{Root cause} & \textbf{Remediation} \\
    \midrule
    F-01 & Grounding failure       & No evidence returned for a fee query    & Evidence-required response; no-fabrication guard \\
    F-02 & Unicode encode error    & Emoji in input printed to cp1252 console & \texttt{stdout} configured with \texttt{errors="replace"} \\
    F-03 & Ticket not found        & Requested ticket ID has no record       & \texttt{NOT\_FOUND} status mapped to clarification outcome \\
    F-04 & Invalid service name    & Identifier not in allowed enum          & Enum validated before tool call; safe handoff on failure \\
    F-05 & Prompt injection        & User attempts to override system prompt & User content treated as untrusted; refusal path taken \\
    \bottomrule
  \end{tabular}
\end{table}

Individual re-test traces are stored as \texttt{week7-T1.json},
\texttt{week7-T3.json}, \texttt{week7-T4.json}, and \texttt{week7-A1.json}
under \texttt{evidence/traces/week7/}.

% ────────────────────────────────────────────────────────────────────────────
\section{Guardrail Implementation}

\subsection{Motivation}

The Week 6 controller contained safety controls embedded inline across its
planning, tool-execution, and observation steps. These controls were effective
but implicit. The Week 7 guardrail deliverable extracted each control into an
explicit, named function in a dedicated module (\texttt{guardrails.py}), making
them independently callable, auditable by name, and testable in isolation.

\subsection{Ten Named Guardrails}

\begin{table}[h]
  \centering
  \caption{Guardrail catalogue --- all ten controls}
  \label{tab:guardrails}
  \small
  \begin{tabular}{@{} C{0.8cm} L{3.5cm} L{4.2cm} L{3.5cm} @{}}
    \toprule
    \textbf{ID} & \textbf{Name} & \textbf{What it blocks} & \textbf{Code reference} \\
    \midrule
    G-01 & Action allow-list        & Unapproved planner actions                 & \texttt{check\_action\_allowlist} \\
    G-02 & Identity injection       & Model-supplied \texttt{customer\_id} or token   & \texttt{check\_identity\_injection} \\
    G-03 & Tool allow-list          & Calls to unapproved tools                  & \texttt{check\_tool\_allowlist} \\
    G-04 & Evidence grounding       & Factual answers without retrieved evidence & \texttt{check\_evidence\_grounding} \\
    G-05 & Escalation approval gate & Escalation without human sign-off          & \texttt{check\_escalation\_approval} \\
    G-06 & Memory-not-authorisation & Memory context auto-approving escalation   & \texttt{check\_memory\_not\_authorization} \\
    G-07 & Credential non-disclosure & Echoing OTPs, PINs, passwords, API keys   & \texttt{check\_credential\_request} \\
    G-08 & Scope boundary           & Funds transfer, loan approval, account change & \texttt{check\_scope} \\
    G-09 & Prompt injection         & Instruction overrides, role substitution   & \texttt{check\_prompt\_injection} \\
    G-10 & Iteration limit          & Loop executing beyond four iterations      & \texttt{check\_iteration\_limit} \\
    \bottomrule
  \end{tabular}
\end{table}

\subsection{Test Evidence}

A dedicated demo script (\texttt{demo\_guardrails.py}) runs 27 test cases,
covering both the ``safe input passes'' and ``unsafe input is blocked'' paths
for every guardrail. The output is written to
\texttt{evidence/traces/week7/guardrail-evidence.json}, recording the guardrail
identifier, test input, trigger result, safety flag, and reason string for each
case. All 27 cases passed with zero failures.

Key correctness properties verified by the tests include: no action or tool
call executes outside the pre-approved sets (G-01, G-03); the model cannot
supply or forge identity fields (G-02); no factual answer is produced when
retrieval returns no evidence, preventing fabricated bank policies (G-04);
prior case memory is read-only context and does not substitute for human
approval of escalation (G-05, G-06); and prompt injection attempts and
out-of-scope financial actions are blocked at input classification (G-08,
G-09).

% ────────────────────────────────────────────────────────────────────────────
\section{Deliverables Summary}

\begin{table}[h]
  \centering
  \caption{Week 7 deliverables and their locations}
  \label{tab:deliverables}
  \small
  \begin{tabular}{@{} L{5.5cm} L{7.0cm} @{}}
    \toprule
    \textbf{Deliverable} & \textbf{File} \\
    \midrule
    30-scenario evaluation set   & \texttt{docs/evaluation/week7/Week\_7\_AI\_Agent\_Evaluation\_Simple\_Test\_Cases.md} \\
    Evaluation harness           & \texttt{evidence/demo/centenary-agent-week6/week7\_evaluation\_harness.py} \\
    Observability trace bundle   & \texttt{evidence/traces/week7/week7-observability-evidence.json} \\
    Failure catalogue            & \texttt{docs/evaluation/week7/Week7\_Failure\_Catalogue.md} \\
    Guardrail module             & \texttt{evidence/demo/centenary-agent-week6/guardrails.py} \\
    Guardrail demo               & \texttt{evidence/demo/centenary-agent-week6/demo\_guardrails.py} \\
    Guardrail evidence           & \texttt{evidence/traces/week7/guardrail-evidence.json} \\
    Guardrail catalogue          & \texttt{docs/evaluation/week7/Week7\_Guardrail\_Catalogue.md} \\
    \bottomrule
  \end{tabular}
\end{table}

% ────────────────────────────────────────────────────────────────────────────
\section{Conclusion}

Week 7 brought the prototype to a state where its behaviour can be evaluated
against a defined scenario set, its internal decisions are fully traceable, its
known failure modes are catalogued and remediated, and its safety controls are
named, implemented, and independently verified. All Week 7 acceptance criteria
were met: all 30 evaluation scenarios are recorded, the five failure cases are
remediated and re-tested as pass, and all 27 guardrail test cases pass with
zero safety violations observed. The agent remains a prototype operating on
synthetic data without production infrastructure, but the evaluation, tracing,
and safety-control foundations are in place for a supervised production
transition.

\vspace{6pt}
\noindent\rule{\textwidth}{0.4pt}
\vspace{2pt}
{\small\noindent\textit{Source code and evidence artefacts: }
\texttt{evidence/demo/centenary-agent-week6/} and
\texttt{evidence/traces/week7/}}

\end{document}
